\documentclass[11pt,a4paper]{article}
\usepackage[utf8]{inputenc}
\usepackage[T1]{fontenc}
\usepackage[margin=2.5cm]{geometry}
\usepackage{booktabs}
\usepackage{tabularx}
\usepackage{enumitem}
\usepackage{parskip}
\usepackage[hidelinks]{hyperref}

\setlist{nosep,leftmargin=1.4em}
\newcommand{\owner}[1]{\textbf{[#1]}}

\title{\vspace{-1.5em}Talking to companies: process, what we send, and who\\[-0.2em]
\large Working note for the ``spreken met bedrijven'' workstream (IOF)}
\author{Robin Wydaeghe \quad (feedback very welcome, esp.\ on contacts: \owner{Luc} \owner{Wout})}
\date{Draft, June 2026}

\begin{document}
\maketitle
\vspace{-1.5em}

\section{Purpose}
The goal of this workstream is to gauge the market, gather a few \textbf{Letters of Intent (LOIs) or
letters of support}, and get \textbf{business-plan validation} from companies, as evidence of market pull
for the IOF ``June'' call (deadline 3 August). This note covers (i) what TechTransfer requires, (ii) what we
send companies and how, (iii) draft emails, and (iv) a first list of contacts with who owns each intro. It is
a draft written to be corrected, the main gap being the contact addresses.

\section{What TechTransfer requires (per A.\ Biondi, 26 May)}
The IP currently sits with UGent, so TechTransfer (TT) runs the confidentiality side. The agreed process:
\begin{itemize}
  \item \textbf{NDA via TT.} The mutual NDA template on the TT SharePoint is the right one. We do not fill in or send it ourselves.
  \item \textbf{Align disclosure first.} Before anything goes out, we agree with TT \emph{what} is disclosed and \emph{in which format}.
  \item \textbf{TT runs the signatures.} We send TT a \textbf{list of people to contact}. TT completes the NDA per recipient, sends it, and collects signatures.
  \item \textbf{Disclose only after signing.} We share information only once \textbf{all} relevant NDAs are signed, and only to the agreed level of detail.
  \item \textbf{Patent gate.} TT asks that \textbf{drafting of the patent application has started} before \emph{any} external sharing, in any format, NDA included. Drafting runs in parallel now and filing is targeted for $\sim$1--15 July.
\end{itemize}

\section{Timeline fit and the deadline tension}
TT estimates the full NDA cycle at \textbf{1--2 months}, which collides with the 3 August IOF deadline. The
sequence resolves this if we start now, because the early steps disclose nothing:
\begin{itemize}
  \item \textbf{Send TT the contact list this week.} Handing over a list of names is not a disclosure, so the NDA clock and the patent drafting both start immediately and run in parallel.
  \item By the time the NDAs come back signed ($\sim$July), the patent draft is under way or filed, so the gate is satisfied exactly when we need to disclose.
  \item \textbf{Keep the target set small.} Fewer NDAs sign faster, and we only need a handful of letters. Do not let the batch wait on one slow legal department.
  \item Pre-agree a single \textbf{teaser} and a single \textbf{NDA deck} with TT, so there is no back-and-forth on scope.
  \item \textbf{Backstop.} A genuinely high-level, non-confidential call needs no NDA at all. We can gauge interest and possibly secure a letter of support without ever showing the technical deck. \owner{Filip}: would such letters of support (rather than signed commercial LOIs) satisfy the IOF jury's ``market pull'' bar? The answer sets how hard we push the NDA track before August.
\end{itemize}

\section{What we send companies, and how}
\paragraph{Two documents.}
\begin{itemize}
  \item \textbf{(A) Non-confidential one-pager (teaser).} The problem, what AEGIS does (per-point absorbed power density on a 3D body in milliseconds, 100\,MHz--300\,GHz), the framing of a fast pre-screen that feeds full-wave validation, two viewer screenshots, and a credentials line. No method internals. Cleared with TT before use.
  \item \textbf{(B) Confidential deck (NDA only, in meeting).} This is the two-concept deck already with Alessandro: pseudo-Brewster compensation and surface confinement, each walked through with its key figure, plus the validation figures and the coherent-MIMO operator. Shown only under a signed NDA and only to the agreed detail. Never attached to an email.
\end{itemize}

\paragraph{How we describe the method (so we stay consistent and protect the patent).}
Lead on what is genuinely and defensibly ours: a \textbf{closed-form, solver-free} evaluation of absorbed
power density over the full body surface in \textbf{milliseconds}, validated analytically (lossy-sphere Mie)
and against published numerical results to within a few percent, positioned as a \textbf{complement} that
pre-screens worst cases ahead of full-wave tools rather than replacing them. Do \emph{not} lead on
``differentiable'' or call it the first differentiable dosimetry method: a group in Split has published
autodiff-based mmWave APD work since 2021, and general differentiable EM is mature, so that framing invites a
prior-art argument we do not need. The transmission-coefficient relationship is likewise textbook. The real
edge is the closed-form speed and the surface formulation, and the pseudo-Brewster and MIMO pieces stay behind
the NDA.

\paragraph{Approach.} A warm introduction beats a cold email every time, so we route through the intros in the
table below wherever they exist. Lead with credibility, keep the tone collaborative and non-commercial (``I
would value your perspective''), make one specific ask, and do not overclaim. The LOI or letter ask comes at
the \emph{end} of a call, once they have seen the demo, never in the first email. No mass emailing and no
public marketing site before the patent is filed.

\section{Draft emails}
\subsection*{Template 1: discovery / letter of support (test lab, regulator, friendly contact)}
\begin{quote}\itshape
Subject: Fast APD pre-compliance method, would value 20 minutes\\[0.3em]
Dear [name],\\[0.3em]
I am Robin Wydaeghe, a final-year PhD with Wout Joseph at UGent (INTEC-WAVES). I work on EM exposure assessment for 5G and 6G, with first-author work in npj Wireless Technology (2026) and two BioEM prizes.\\[0.3em]
Over the last months I have built a method that computes absorbed power density on a 3D body surface in milliseconds, across 100\,MHz to 300\,GHz, validated against published numerical results to within a few percent. It is meant as a fast pre-screening layer ahead of full-wave tools, not a replacement.\\[0.3em]
I am trying to understand whether this is useful to teams like yours, and would value 20 to 30 minutes to show a short demo and hear how [APD pre-compliance / site compliance] works on your side. A one-page summary is attached.\\[0.3em]
Would a short call in the next two weeks work? Best, Robin
\end{quote}

\subsection*{Template 2: ZMT / Sim4Life (peer, collaboration, not a sales pitch)}
\begin{quote}\itshape
Subject: Closed-form APD method, possible matched validation against Sim4Life\\[0.3em]
Dear [Sven / name],\\[0.3em]
I am Robin Wydaeghe, PhD with Wout Joseph at UGent. You may know our group through the IEC TC\,106 and 63195 work. I previously placed in a Sim4Life competition with a large automation pipeline (GOLIAT), so I work inside your ecosystem rather than against it.\\[0.3em]
I have developed a closed-form surface formulation for absorbed power density at mmWave that runs in milliseconds and matches published numerical results to a few percent. I see it as a fast pre-screening complement to Sim4Life: find the worst-case beam, grip and position scenarios in seconds, then run those rigorously in Sim4Life.\\[0.3em]
Before anything else I would like to do a clean matched validation on a standard 63195 reference scenario. If the figures hold up, I think there is a natural joint contribution, possibly co-presented at BioEM. Would you be open to a short call? Best, Robin
\end{quote}

\section{Target contacts (first attempt, please correct)}
Addresses are to be sourced or confirmed by the owner in the last column. I have not guessed personal email
addresses. The two groups matter: the \textbf{fast track} is the friendly, low-stakes end where a contact can
give a letter cheaply before August, and the \textbf{relationship track} is the patient, post-patent set that
should \emph{not} be rushed for an August LOI. Per Wout, Sven K\"uhn is the spin-off and partnership lane,
while Azadeh Peyman is kept for the competition and jury side, so those two stay in separate lanes.

\paragraph{Fast track (IOF letters before 3 August).}
\begin{tabularx}{\textwidth}{@{}p{3.0cm} X p{3.9cm}@{}}
\toprule
\textbf{Organisation / person} & \textbf{Why it matters} & \textbf{How we reach them (owner)} \\
\midrule
Ericsson Research --- D.\ Colombi, C.\ T\"ornevik & Base-station and massive-MIMO actual-max EMF compliance; strongest LOI lead & Wout $\to$ Luc Martens $\to$ both, 2-hop \owner{Wout}\,\owner{Luc} \\
Test lab --- Eurofins E\&E (DE), CETECOM (Essen), Verkotan (Oulu) & Device pre-compliance; short decision cycle, a lab lead can sign alone & Public pages; K\"uhn interfaces with all three \owner{Robin} \\
Regulator --- BIPT (Belgium) & Local; letter of relevance and standards leverage & Wout knows the BIPT EMF engineer \owner{Wout} \\
Regulator --- ARPANSA (Australia) & National authority; reachable via Carolina's grant & Carolina \owner{Robin} \\
Operator --- Orange Innovation (D.-T.\ Phan-Huy) & EMF-aware / RIS research plus operator market & via Margot Deruyck (INTEC-WAVES), 2nd-degree \owner{Robin} \\
Operator --- Proximus / Telenet (BE) & Local operators; site-compliance use case & \owner{Luc}\,\owner{Wout} \\
Key stakeholders named in the IDF & Already identified as relevant; extract and confirm addresses & \owner{Robin} + \owner{Luc}\,\owner{Wout} \\
\bottomrule
\end{tabularx}

\paragraph{Relationship track (post-patent, do not rush for an August letter).}
\begin{tabularx}{\textwidth}{@{}p{3.0cm} X p{3.9cm}@{}}
\toprule
\textbf{Organisation / person} & \textbf{Why it matters} & \textbf{How we reach them (owner)} \\
\midrule
ZMT / Sim4Life --- Sven K\"uhn (+ N.\ Kuster, E.\ Neufeld) & Sim4Life ecosystem; matched validation $\to$ co-present $\to$ integrate. Wout endorsed Sven for the spin-off lane & Wout, via IT'IS / IEC TC\,106 \owner{Wout} \\
SPEAG / IT'IS (Z\"urich) & Measurement and tissue-data orbit; standards reach & \owner{Wout} \\
Nokia & RAN-vendor compliance & via Jafar Keshvari (IEEE ICES Chair), 2nd-degree through Luc and Mike Wood; ask who owns EMF at Nokia \owner{Luc}\,\owner{Wout} \\
\bottomrule
\end{tabularx}

\section{What I need from you}
\begin{itemize}
  \item \owner{Wout}\,\owner{Luc}: corrections and missing addresses for the tables above. This is the main gap. In particular the BIPT engineer's name (Wout), the Nokia EMF owner via Keshvari, and a green light to use the Ericsson intro through Luc.
  \item \owner{Wout}: availability (which days work and which do not) for the company meetings you want to join.
  \item \owner{Filip}: does the IOF jury accept letters of support and documented interest, or does it need signed commercial LOIs? This sets how hard we push the NDA track before 3 August.
  \item \owner{TT (Alessandro)}: agree the disclosure scope, one cleared teaser (A) and one NDA-only deck (B), and confirm the per-recipient NDA timing once the contact list is in.
\end{itemize}

\end{document}
